% Reviewed FF.36--FF.37, integrated after root and geometric_routes audits. Do not input this duplicate.
\begin{claim}{FF.36}{Proved}\label{clm:FFabelianHT}
Let \(M/\mathbb Q_p\) be finite, and let \(S\)
be a Hodge--Tate representation of \(G_M\) whose
inertia image is abelian. The inertia representation
on \(S\) is semisimple. In particular there is a
canonical decomposition
\[
 S=S^{I_M}\oplus
       \sum_{g\in I_M}(g-1)S.                         \tag{FF.36.1}
\]
Its projection onto the first summand commutes
with \(G_M\). For every inertia-stable subspace
\(N\), that projection sends \(N\) to
\(N\cap S^{I_M}\), and taking inertia invariants
is exact on \(0\to N\to S\to S/N\to0\).
All assertions hold after any finite extension
of \(M\).
\end{claim}
\begin{proof}
Write \(\mathfrak g\) for the Lie algebra of
the compact inertia image. It is abelian.
Let \(B\subset\operatorname{End}_{\mathbb Q_p}(S)\)
be the unital commutative algebra generated by
\(\mathfrak g\). It is a finite-dimensional
commutative algebra over the characteristic-zero
field \(\mathbb Q_p\).

We recall the elementary algebra fact used
here. Such an algebra has a unique maximal
\'etale subalgebra \(B_{\rm ss}\), and
\[
           B=B_{\rm ss}\oplus\operatorname{rad}(B)
                         \quad\text{as vector spaces}.
                                                        \tag{FF.36.2}
\]
One proof uses the additive and multiplicative
Jordan decompositions of commuting matrices:
their semisimple parts belong to the algebra,
are polynomials in the original matrices, and
add and multiply as the original matrices do
modulo the nilpotent radical. These parts
give the unique section of the finite separable
algebra \(B/\operatorname{rad}(B)\). Equivalently,
in each local Artin factor, its separable residue
field has a unique lift through the nilpotent
ideal. After extension to \(\mathbb C_p\),
each factor has the form
\(\mathbb C_p\oplus\mathfrak m\), with
\(\mathfrak m\) nilpotent. An element acting
semisimply on the faithful module has zero
\(\mathfrak m\)-component. Thus the semisimple
elements of \(B\otimes\mathbb C_p\) are exactly
\(B_{\rm ss}\otimes\mathbb C_p\).

The Sen operator \(\Theta\) of \(S\) is
semisimple, because \(S\) is Hodge--Tate
\citep[Section~2.7]{FF_Fontaine2004}.
By \citet[Theorem~3.1, \(d=0\)]{FF_Ohkubo2014},
\(\Theta\in\mathfrak g\otimes\mathbb C_p\),
and \(\mathfrak g\) is the smallest
\(\mathbb Q_p\)-subspace of the endomorphism
space with this property. The preceding
algebra fact puts \(\Theta\) in
\(B_{\rm ss}\otimes\mathbb C_p\). Minimality
therefore gives \(\mathfrak g\subset B_{\rm ss}\).
In particular every element of \(\mathfrak g\)
acts semisimply, and these elements commute.

On a sufficiently small open subgroup of
the compact inertia image, matrix logarithm
has values in \(\mathfrak g\), and matrix
exponential is its inverse. Hence every
element of this open subgroup is semisimple.
An arbitrary inertia matrix has a positive
integer power in this subgroup. A matrix
over a characteristic-zero field whose
positive power is invertible and semisimple
is itself semisimple: if that power has
distinct eigenvalues \(\alpha_j\), the
matrix is annihilated by the square-free
polynomial \(\prod_j(T^n-\alpha_j)\).
All inertia matrices commute, so they are
simultaneously diagonalizable over a finite
extension of \(\mathbb Q_p\). Their
representation is semisimple over
\(\mathbb Q_p\) as well.

In this simultaneous decomposition,
\(S^{I_M}\) is the trivial-character part,
and the sum in (FF.36.1) is exactly the
sum of the nontrivial-character parts.
This description is independent of a
coefficient extension. It also proves
the assertion about every stable \(N\)
and exactness of invariants. Conjugation
by \(G_M\) preserves inertia and its
trivial-character part, as well as the
explicit second summand in (FF.36.1).
The projection therefore commutes with
\(G_M\). Restricting to any open local
Galois subgroup retains both the
Hodge--Tate and abelian-image hypotheses.
\end{proof}

\begin{claim}{FF.37}{Proved}\label{clm:FFmaxabelianreconstruct}
Let \(K_p=\operatorname{Frac}(\mathbb Z_{(p)}^{sh})\)
with the chosen embedding in \(\overline{\mathbb Q}\),
and let \(K_p^{\rm ab}\) be its maximal
abelian extension in that algebraic closure.
Every finite extension of \(K_p^{\rm ab}\),
and every subfield of such an extension,
satisfies \'etale reconstruction.
The extension \(K_p^{\rm ab}/U_p\) is infinite,
so this is a proper enlargement of
Claim~\ref{clm:FFcyclotomicreconstruct} at the
level of the specified henselian fields.
\end{claim}
\begin{proof}
Put \(I=I_{\mathbb Q_p}=G_{K_p}\), and
let \(J=\overline{[I,I]}\) be its closed
commutator subgroup. Thus
\(G_{K_p^{\rm ab}}=J\).
The tame quotient of inertia is abelian,
so \(J\subset P_{\mathbb Q_p}\).
It contains the commutator subgroup of
the free pro-\(p\) group
\(P_{\mathbb Q_p}\), and is infinite:
two free generators have a nonidentity
commutator, of infinite order by
torsion-freeness. Every open subgroup
of \(J\) is consequently infinite.
This uses the same verified free wild
inertia theorem as Claim~\ref{clm:FFcyclotomictest}.

Both \(T_p/K_p\) and the cyclotomic
extension \(K_p(\mu_{p^\infty})/K_p\)
are abelian. Hence \(U_p\subset K_p^{\rm ab}\).
In particular these fields contain
\(\mu_\infty\), have value group
\(\mathbb Q\), and have residue field
\(\overline{\mathbb F}_p\).

\emph{1. The enlargement is infinite.}
Set \(L=K_p(\zeta_p)\), including
\(L=K_2\) when \(p=2\), and let
\(\pi=\zeta_p-1\). Then \(\pi\) is
a uniformizer of \(L\), and
\(\Delta=\operatorname{Gal}(L/K_p)\)
has order \(p-1\). For the principal
unit filtration \(U^i=1+\pi^i\mathcal O_L\),
\[
                    U^1/U^2\simeq
                 \overline{\mathbb F}_p
                \quad\text{as }\mathbb F_p\text{-spaces}.
                                                        \tag{FF.37.1}
\]
The \(\Delta\)-action on this quotient
is the mod-\(p\) cyclotomic character
\(\bar\chi_p\): its residue action is
trivial, and
\(\sigma(\pi)\equiv\bar\chi_p(\sigma)\pi
\pmod{\pi^2}\).
If \(u^p\in U^1\), then \(u\) is a
unit of residue one, because the residue
field has characteristic \(p\). Writing
\(u=1+\pi c\), the binomial expansion
and \(v_L(p)=p-1\) show that
\(u^p\in U^2\). Therefore (FF.37.1)
is a quotient of a subspace of
\(L^\times/(L^\times)^p\).
Since \(|\Delta|\) is prime to \(p\),
taking its character parts is exact.
The \(\bar\chi_p\)-part of this Kummer
space is infinite-dimensional over
\(\mathbb F_p\).

Kummer theory and the twist by
\(\mu_p^{-1}\) identify that part
with \(H^1(G_L,\mathbb F_p)^\Delta\).
Inflation--restriction for \(\Delta\)
identifies the latter with
\(H^1(I,\mathbb F_p)\), because
the positive-degree cohomology of this
prime-to-\(p\) finite group is zero.
Thus \(I^{\rm ab}\) has an
infinite-dimensional elementary abelian
pro-\(p\) quotient. By contrast
\(\operatorname{Gal}(U_p/K_p)\)
is a closed subgroup of a product
of a pro-prime-to-\(p\) procyclic
group and a subgroup of
\(\mathbb Z_p^\times\); its pro-\(p\)
Sylow embeds in the second factor.
Its elementary abelian pro-\(p\)
quotient is finite-dimensional.
A finite kernel could increase that
dimension by only a finite amount.
Hence \(K_p^{\rm ab}/U_p\) cannot be finite.

\emph{2. Exact cohomological inputs for all finite extensions.}
Let \(K/K_p^{\rm ab}\) be finite.
It has the form \(L_0K_p^{\rm ab}\)
for a finite extension \(L_0/K_p\).
Every specified element of \(K\)
can be included in a larger such
\(L_0\). The extension \(K/L_0\)
is Galois abelian and has infinite
cyclotomic image. The abelian cocycle
calculation (FF.33.2), together with
the exact Kummer kernel over \(L_0\),
therefore proves
\[
 \ker\bigl(K^\times\to H^1(G_K,\mathbb Z_p(1))\bigr)
 =\ker\bigl(K^\times\to H^1(G_K,\mathbb Q_p(1))\bigr)
 =\mu_\infty.
                                                        \tag{FF.37.2}
\]
Compatible roots may be chosen because
all roots of unity are in \(K\).
This is the proof of
Claim~\ref{clm:FFcyclotomiccohomology}
with another abelian extension in
place of \(U_p\); all its hypotheses
have just been checked.

Choose a number field \(D_0\subset K\)
containing any prescribed finite data
and generators of \(K/K_p^{\rm ab}\),
and put \(M_0=(D_0)_v\). Inside the
fixed local Galois group the absolute
group of \(K\) is
\[
            H_0=J\cap G_{M_0}
                         \triangleleft G_{M_0}.
                                                        \tag{FF.37.3}
\]
The quotient \(I_{M_0}/H_0\) is
abelian, as a subgroup of \(I/J\),
and its cyclotomic image is infinite.
On this abelian quotient, (FF.33.2)
gives \(H^1(I_{M_0}/H_0,\mathbb Q_p(1))=0\).
Its coefficient invariants are zero.
Inflation--restriction for inertia
in \(G_{M_0}/H_0\) then also gives
\(H^1(G_{M_0}/H_0,\mathbb Q_p(1))=0\).
Consequently both restriction maps
\[
 H^1(I_{M_0},\mathbb Q_p(1))\longrightarrow H^1(H_0,\mathbb Q_p(1)),
 \qquad
 H^1(G_{M_0},\mathbb Q_p(1))\longrightarrow H^1(H_0,\mathbb Q_p(1))
                                                        \tag{FF.37.4}
\]
are injective. Every finite extension
\(E/K\) equals \(D_1K_p^{\rm ab}\)
for a finite number field \(D_1/D_0\),
and has \(H_1=J\cap G_{M_1}\) with
\(M_1=(D_1)_v\). The same proof gives
(FF.37.2)--(FF.37.4) at every such stage.

\emph{3. The invisible subspace.}
Use the regular pointed source \(X\),
its generating generalized Albanese
\(j:X\to A\), the actual point function
\(b\), and its geometric functional
\(\lambda\) as in
Claim~\ref{clm:FFcyclotomicreconstruct}.
Include the data in \(D_0\), and set
\[
                   S=\operatorname{Hom}_{H_0}
                                (V_pA,\mathbb Q_p(1)).
                                                        \tag{FF.37.5}
\]
It is a \(G_{M_0}\)-stable Hodge--Tate
space, by the exact comparison and
subobject results verified in
Claim~\ref{clm:FFcyclotomicsplitting}.
Its inertia action factors through
the abelian quotient in (FF.37.3).
Claim~\ref{clm:FFabelianHT} thus applies.

Let \(N\subset S\) consist of those
functionals whose pushforward of
\(\kappa_A(a)\) is zero for every
\(a\in A(L)\) over every finite
extension \(L/K\). This is a linear
subspace. It is \(G_{M_0}\)-stable:
conjugation sends such an \(L\) to
another finite extension of \(K\),
and sends the point and its Kummer
cocycle to their conjugates. This
uses normality of \(H_0\), and
includes the coefficient action.

The same \(N\) is the common kernel
of all evaluations at \(j(x)\).
To see the nontrivial inclusion, recall
that finitely many signed sums of
\(j(X_{\bar K})\) dominate \(A_{\bar K}\).
Every point of \(A(\bar K)\) is a
difference of points in this dominant
image, since two nonempty opens related
by a translation intersect. Thus it
is a finite signed sum of \(j(x_i)\),
after a finite field extension.
A nonzero rational Kummer evaluation
cannot become zero under a finite
restriction, by restriction--corestriction.
Consequently vanishing on all \(j(x)\)
implies vanishing on every \(a\).
Finite-dimensionality of \(S/N\)
now supplies finitely many such
evaluations with common kernel
exactly \(N\), over one finite
extension \(E/K\).

For every finite local model \(M_1\)
obtained as in Step~2, we have
\[
                        N\cap S^{I_{M_1}}=0.
                                                        \tag{FF.37.6}
\]
Indeed, a nonzero inertia-invariant
functional is detected over a finite
extension of the corresponding strict
henselization by
Claim~\ref{clm:FFpdetector}. Its
nonzero class survives adjoining
\(K_p^{\rm ab}\) by (FF.37.4),
so it cannot belong to \(N\).

\emph{4. The canonical invariant lift.}
Include the finitely many detecting
points and their actual values in
\(D_1\subset E\), so that
\(E=D_1K_p^{\rm ab}\) and
\(H_1=J\cap G_{M_1}\) is normal
in \(G_{M_1}\). The actual-value
cohomology identity, followed by
conjugation, gives
\[
                  g\lambda-\lambda\in N
                         \quad(g\in G_{M_1}).
                                                        \tag{FF.37.7}
\]
Here the selected point and value
cocycles are defined over \(M_1\),
and their restricted classes are
conjugation-invariant. Their common
evaluation kernel is exactly \(N\).
No descent of the topos map is
being assumed.

Take the canonical projection
\(\Pi:S\to S^{I_{M_1}}\) from
Claim~\ref{clm:FFabelianHT}, and put
\(\lambda_0=\Pi\lambda\).
Equation (FF.37.7) makes the image
of \(\lambda\) in \(S/N\)
inertia-invariant. Exactness of
invariants and (FF.37.6) imply
\[
               \lambda-\lambda_0\in N,
           \qquad \lambda_0\in
             \operatorname{Hom}_{G_{M_1}}(V_pA,\mathbb Q_p(1)).
                                                        \tag{FF.37.8}
\]
For the second assertion, \(\Pi\)
commutes with \(G_{M_1}\), and
\(\Pi N=N\cap S^{I_{M_1}}=0\).
Thus applying \(\Pi\) to
(FF.37.7) kills every difference.
This is the unique inertia-invariant
lift of the class modulo \(N\).
It is not asserted to have been
\(I_{M_0}\)-invariant before the
finite actual-value descent.

After splitting the torus of \(A\)
over a further finite model, restriction
of a functional to that torus has
inertia-trivial target
\(\operatorname{Hom}(V_pT,\mathbb Q_p(1))\).
Every element of \(N\) has zero
restriction: an inertia-equivariant
map to a trivial representation
factors through \(\Pi\), and
\(\Pi N=0\). The original integral
boundary residues of \(\lambda\)
are therefore unchanged by replacement
with \(\lambda_0\).

\emph{5. Reconstruction from the retained point values.}
Construct the normalized rational
class \(\xi_0\) from \(\lambda_0\)
using the section at \(x_0\),
exactly as in Step~1 of
Claim~\ref{clm:FFcyclotomicreconstruct}.
Since \(\lambda-\lambda_0\in N\),
its evaluation at every finite-extension
point is \(\kappa_p(b(x)/b_0)\).
Its residues are the integral residues
of the original class, by Step~4.
All the subsequent hypotheses in
that reconstruction proof are now
available, as follows.

First, (FF.37.4) lifts an equality
of actual evaluations from \(H_2\)
to the finite local model \(G_{M_2}\).
The two nearby punctured-line values
therefore force each boundary residue
into \(\mathbb Q\cap\mathbb Z_p\),
using the ordinary finite-local-field
valuation. Norm transfer and transverse
curves give the same assertion in
every dimension. Second, the invariant
rational Picard 1-motive lift can
be multiplied to have integral Tate
coordinates and an integral lattice
coordinate. The compact local Picard
group then makes a further multiple
of that boundary divisor principal,
and its invariant zero-residue Tate
part is zero. Thus a unit \(h\)
and integer \(d>0\) satisfy
\[
                    b(x)^d=\zeta_xh(x),
                         \qquad\zeta_x\in\mu_\infty\subset K,
                                                        \tag{FF.37.9}
\]
by (FF.37.2). This is an equality
of point values, not an identification
of the original integral topos class
with a unit class.

Third, every finite extension of
\(K\) has infinite absolute Galois
group, since its group is open in
\(J\). Hence the primitive-element
and finite-map argument in Step~4
of Claim~\ref{clm:FFcyclotomicreconstruct}
supplies closed points of arbitrarily
large degree on every nonempty
curve open. The \(T,T-1\) resultant
applied to (FF.37.9) therefore makes
the retained nonzero functional the
geometric Kummer functional of a unit.
If it is zero, every actual value
lies in \(b_0\mu_\infty\subset K\),
and the \(T-\beta\) argument forces
the point function to be constant.
The \(T-a\) comparisons remove the
remaining root-of-unity scalars.

These comparisons are compatible
with the canonical projections just
used. To verify this without making
\(N\) explicit, a pullback on
Albanese Tate modules carries an
inertia-invariant realizing functional
to an inertia-invariant one, and
carries an invisible functional to
an invisible one: point Kummer
classes pull back by actual points.
Uniqueness of the invariant lift
in (FF.37.8) gives the required
compatibility, after a common finite
model extension. The same applies
to restrictions to nonempty opens
and to the curves used for testing.
Thus every step of the actual
point-function and scalar-removal
proof is preserved.

It follows that the pinned map's
geometric point function is realized
by a unit over a finite extension.
As in Claim~\ref{clm:FFstrictreconstruct},
equality on geometric points makes
that unit Galois-invariant, so it
descends to the original \(K\).
The point-function criterion
\citep[Theorem~4.18(5)]{FF_CHW2024}
proves \(\mathcal R(K)\).
Since \(K\) is algebraic over
\(\mathbb Q\), it is algebraic
over each of its subfields.
Claim~\ref{clm:FFdescent} proves
the final subfield assertion.
\end{proof}
